%% ma: L11-B02-0004
%% lop: 11
%% chuong: 1
%% bai: 2
%% dang: 11.2.7
%% loai: TL
%% muc: VD
%% dap_an: "-"
%% nguon: "de-cuong-giua-ky-1-toan-11-nam-2026-2027-truong-thpt-hoang-van-thu-ha-noi-2-8.pdf (M08, Đề cương giữa kì 1 Toán 11 THPT Hoàng Văn Thụ, Hà Nội), Phần 4 Bài 1 ý 1"
%% kien_thuc: "Bài 1 (giá trị lượng giác của góc lượng giác), Bài 2 (công thức lượng giác)"
%% trang_thai: cho-duyet
%% ly_do: "Dạng mới đề xuất 11.2.7"
\begin{ex}
Cho $\cos\alpha=-\dfrac{5}{13}$ và $\dfrac{\pi}{2}<\alpha<\pi$. Tính $\sin\alpha$; $\cos2\alpha$; $\sin\left(\alpha-\dfrac{\pi}{6}\right)$; $\tan2\alpha$; $\cos\dfrac{\alpha}{2}$.
\loigiai{Vì $\dfrac\pi2<\alpha<\pi$ nên $\sin\alpha>0$: $\sin\alpha=\sqrt{1-\dfrac{25}{169}}=\dfrac{12}{13}$.

$\cos2\alpha=2\cos^2\alpha-1=\dfrac{50}{169}-1=-\dfrac{119}{169}$.

$\sin\left(\alpha-\dfrac\pi6\right)=\sin\alpha\cdot\dfrac{\sqrt3}2-\cos\alpha\cdot\dfrac12=\dfrac{12\sqrt3+5}{26}$.

$\sin2\alpha=2\sin\alpha\cos\alpha=-\dfrac{120}{169}$ nên $\tan2\alpha=\dfrac{\sin2\alpha}{\cos2\alpha}=\dfrac{120}{119}$.

$\dfrac\pi4<\dfrac\alpha2<\dfrac\pi2$ nên $\cos\dfrac\alpha2>0$: $\cos^2\dfrac\alpha2=\dfrac{1+\cos\alpha}2=\dfrac4{13}$, $\cos\dfrac\alpha2=\dfrac{2\sqrt{13}}{13}$.}
\end{ex}
